% ============================================================
%  Stage_M2_IMD_modelisation_csp.tex
%  Chapitre 1 du rapport de stage M2 IMD --- AMU
%  Modelisation par graphes et CSP du probleme de generation
%  de molecules non-benzenoides.
% ============================================================
\documentclass[12pt,a4paper]{article}

\usepackage[T1]{fontenc}
\usepackage[utf8]{inputenc}
\usepackage[french]{babel}
\usepackage{geometry}
\usepackage{amsmath,amssymb,amsthm}
\usepackage{mathtools}
\usepackage{array}
\usepackage{booktabs}
\usepackage{colortbl}
\usepackage{enumitem}
\usepackage{xcolor}
\usepackage{fancyhdr}
\usepackage{titlesec}
\usepackage{tcolorbox}
\usepackage{listings}
\usepackage{parskip}
\usepackage{longtable}
\usepackage{tikz}
\usepackage{algorithm}
\usepackage{algpseudocode}
\usepackage{hyperref}

\usetikzlibrary{calc,positioning,arrows.meta,shapes.geometric,backgrounds}
\tcbuselibrary{skins,breakable}

\hypersetup{
  colorlinks=true,
  linkcolor=blue!50!black,
  citecolor=blue!50!black,
  urlcolor=blue!50!black,
}

\geometry{
  top=2.5cm, bottom=2.5cm,
  left=3.0cm, right=2.5cm,
  headheight=15pt
}

\definecolor{bleuTitre}{RGB}{31,56,100}
\definecolor{bleuSection}{RGB}{31,92,153}
\definecolor{bleuSub}{RGB}{46,117,182}
\definecolor{grisNote}{RGB}{245,244,240}
\definecolor{codebg}{RGB}{248,246,240}
\definecolor{codecomment}{RGB}{100,100,90}
\definecolor{vertOk}{RGB}{30,120,50}
\definecolor{rougeAlert}{RGB}{180,40,20}
\definecolor{grisLigne}{RGB}{235,235,230}
\definecolor{jaunePent}{RGB}{254,215,170}
\definecolor{grisHex}{RGB}{203,213,225}
\definecolor{indigoHept}{RGB}{165,180,252}

\pagestyle{fancy}
\fancyhf{}
\renewcommand{\headrulewidth}{0.4pt}
\fancyhead[L]{\small\textit{Chapitre 1 -- Mod\'elisation CSP}}
\fancyhead[R]{\small\textit{Stage M2 IMD}}
\fancyfoot[C]{\thepage}

\titleformat{\section}
  {\large\bfseries\color{bleuTitre}}
  {\thesection.}{0.6em}{}[\vspace{-4pt}\rule{\linewidth}{0.6pt}]
\titleformat{\subsection}
  {\normalsize\bfseries\color{bleuSection}}
  {\thesubsection}{0.6em}{}
\titleformat{\subsubsection}
  {\normalsize\itshape\bfseries}
  {\thesubsubsection}{0.6em}{}

\theoremstyle{definition}
\newtheorem{definition}{D\'efinition}[section]
\newtheorem{propriete}[definition]{Propri\'et\'e}
\newtheorem{contrainte}[definition]{Contrainte}
\newtheorem{theoreme}[definition]{Th\'eor\`eme}
\newtheorem{lemme}[definition]{Lemme}
\newtheorem{corollaire}[definition]{Corollaire}
\theoremstyle{remark}
\newtheorem{remarque}[definition]{Remarque}
\newtheorem{exemple}[definition]{Exemple}

\newtcolorbox{notebox}[1][]{
  enhanced, breakable,
  colback=grisNote,
  colframe=bleuSection,
  leftrule=4pt, toprule=0.4pt, bottomrule=0.4pt, rightrule=0.4pt,
  arc=2pt,
  fonttitle=\bfseries\small\color{bleuSection},
  title={#1},
  before upper={\setlength{\parskip}{4pt}},
}

\lstdefinestyle{benzai}{
  backgroundcolor=\color{codebg},
  basicstyle=\ttfamily\small,
  breaklines=true,
  frame=single,
  framesep=6pt,
  rulecolor=\color{bleuSection!40},
  xleftmargin=12pt,
  xrightmargin=6pt,
  commentstyle=\color{codecomment}\itshape,
  morecomment=[l]{\#},
  showstringspaces=false,
}

\lstdefinestyle{pseudo}{
  backgroundcolor=\color{codebg!60},
  basicstyle=\ttfamily\small,
  breaklines=true,
  frame=leftline,
  framesep=4pt,
  rulecolor=\color{bleuSection!60},
  xleftmargin=10pt,
  commentstyle=\color{codecomment}\itshape,
  showstringspaces=false,
  mathescape=true,
}

% ── Macros ─────────────────────────────────────────────────
\newcommand{\B}{\mathcal{B}}
\newcommand{\GD}{G_{\!D}}
\newcommand{\GB}{G_{\!\B}}
\newcommand{\VD}{V_{\!D}}
\newcommand{\ED}{E_{\!D}}
\newcommand{\xv}{x_v}
\newcommand{\degv}{\deg(v)}
\newcommand{\Tn}{\mathcal{T}(n)}
\newcommand{\ncinq}{n_5}
\newcommand{\nsept}{n_7}
\newcommand{\nhex}{n_6}
\newcommand{\posIn}[2]{\hat{p}_{#1}(#2)}
\newcommand{\pat}[1]{\mathtt{#1}}
\newcommand{\Aut}{\mathrm{Aut}}
\newcommand{\dom}{\mathrm{dom}}
\newcommand{\NN}{\mathcal{N}}
\newcommand{\bigO}{\mathcal{O}}

\renewcommand{\algorithmicrequire}{\textbf{Entr\'ee :}}
\renewcommand{\algorithmicensure}{\textbf{Sortie :}}

% ────────────────────────────────────────────────────────────
\begin{document}

% ============================================================
%  Page de titre du chapitre
% ============================================================
\begin{center}
  {\LARGE\bfseries\color{bleuTitre}
    Chapitre 1\\[6pt]
    Mod\'elisation par graphes et CSP\\[4pt]
    du probl\`eme de g\'en\'eration de mol\'ecules non-benz\'eno\"ides}\\[18pt]
  {\large\itshape
    Approche th\'eorique, alg\'ebrique et algorithmique}\\[14pt]
  {\small Rapport de stage --- Master 2 Informatique et Math\'ematiques Discr\`etes}\\[2pt]
  {\small Aix-Marseille Universit\'e --- \today}
\end{center}

\vspace{8pt}
\rule{\linewidth}{1pt}
\vspace{4pt}

\tableofcontents

\vspace{10pt}
\rule{\linewidth}{0.4pt}
\vspace{8pt}

\begin{notebox}[R\'esum\'e du chapitre]
Nous d\'ecrivons dans ce chapitre la mod\'elisation th\'eorique du probl\`eme
combinatoire qui sous-tend l'ensemble du stage : \'enum\'erer toutes les
mol\'ecules \emph{non-benz\'eno\"ides} obtenues en substituant certains cycles
hexagonaux d'un benz\'eno\"ide par des cycles pentagonaux ou heptagonaux,
tout en conservant le nombre total de carbones. Notre approche repose sur
trois objets math\'ematiques compl\'ementaires : le \emph{graphe de
carbone} pour repr\'esenter l'entr\'ee, le \emph{graphe dual} pour
formuler le probl\`eme, et un \emph{probl\`eme de satisfaction de
contraintes} (CSP) pour le r\'esoudre. Nous \'etablissons les propri\'et\'es
structurelles du graphe dual qui permettent un pr\'e-traitement efficace,
formulons six contraintes (cinq fixes plus une optionnelle) avec leurs
justifications, puis \'etudions les algorithmes de parcours et d'isomorphisme
n\'ecessaires \`a la reconstruction tridimensionnelle des structures
\'enum\'er\'ees. Une attention particuli\`ere est port\'ee aux liens entre la
\textbf{th\'eorie des graphes planaires} (th\'eor\`eme de Whitney,
2-connexit\'e) et la \textbf{combinatoire des cycles d'un benz\'eno\"ide}.
\end{notebox}

% ════════════════════════════════════════════════════════════
\section{Introduction et probl\'ematique}
% ════════════════════════════════════════════════════════════

\subsection{Cadre chimique et motivation}

Un \textbf{hydrocarbure benz\'eno\"ide} $\B$ est une mol\'ecule plane
constitu\'ee de cycles hexagonaux fusionn\'es, chaque sommet
repr\'esentant un atome de carbone et chaque ar\^ete une liaison covalente
C--C. Cette famille comprend de nombreux \emph{hydrocarbures aromatiques
polycycliques} (HAP) tels que le napht\'al\`ene, l'anthrac\`ene ou le
coron\`ene. Les benz\'eno\"ides poss\`edent une structure r\'eguli\`ere :
les angles internes valent~$120^\circ$ et toutes les liaisons C--C sont
sensiblement de m\^eme longueur ($\approx 1{,}40$~\AA{} en environnement
aromatique).

Nous nous int\'eressons \`a une famille \'elargie de mol\'ecules :
celles obtenues en \emph{rempla\c{c}ant} certains cycles hexagonaux par des
cycles \textbf{pentagonaux} (5~carbones) ou \textbf{heptagonaux}
(7~carbones), \`a nombre total d'atomes de carbone constant. Ces structures
\og non-benz\'eno\"ides \fg{} pr\'esentent un int\'er\'et chimique majeur :
\begin{itemize}
  \item l'introduction de cycles \`a 5~ou~7~atomes brise la planarit\'e
        \emph{a priori} de la mol\'ecule (\og courbure de Gauss \fg{}
        positive en pr\'esence de pentagones, n\'egative pour les
        heptagones) ;
  \item lorsque pentagones et heptagones apparaissent en paires
        adjacentes (\emph{motif azul\'enique}), leurs courbures se
        compensent et la mol\'ecule peut retrouver une g\'eom\'etrie plane,
        avec des propri\'et\'es \'electroniques distinctes ;
  \item certaines de ces structures sont des \textbf{biradicaux} stables
        (le c\'el\`ebre \emph{triangul\`ene} de Clar) ou des candidats
        pour de nouveaux mat\'eriaux semi-conducteurs organiques.
\end{itemize}

\subsection{Formulation du probl\`eme}

\noindent\textbf{Probl\`eme.} \'Etant donn\'e un benz\'eno\"ide $\B$
\`a $h$~hexagones, \'enum\'erer toutes les substitutions admissibles
$(x_1, x_2, \ldots, x_h) \in \{5, 6, 7\}^h$ des tailles de ses cycles, telles
que :
\begin{enumerate}[label=(\roman*),leftmargin=*]
  \item le nombre total d'atomes de carbone soit conserv\'e ;
  \item la mol\'ecule r\'esultante reste r\'ealisable g\'eom\'etriquement ;
  \item chaque configuration locale (cycle central + voisins) appartienne
        \`a un catalogue de configurations \og physiquement plausibles \fg{}
        pr\'e-calcul\'e par dynamique mol\'eculaire ;
  \item les solutions \og structurellement \'equivalentes \fg{} (li\'ees
        par une sym\'etrie du benz\'eno\"ide d'entr\'ee) ne soient
        \'enum\'er\'ees qu'une fois.
\end{enumerate}

\subsection{Approche : CSP sur le graphe dual}

Une \'enum\'eration na\"ive sur $\{5,6,7\}^h$ pose deux probl\`emes :
\begin{enumerate}[leftmargin=*]
  \item la taille de l'espace ($3^h$) cro\^it exponentiellement avec
        le nombre d'hexagones : $3^{20} \approx 3{,}5 \cdot 10^9$
        configurations \`a tester ;
  \item la majorit\'e des configurations sont \emph{trivialement
        invalides} (non-conservation du carbone, voisinages topologiquement
        infaisables, doublons par sym\'etrie).
\end{enumerate}

Nous adoptons une approche par \textbf{satisfaction de contraintes}
(CSP, \emph{Constraint Satisfaction Problem}) formul\'ee sur le
\emph{graphe dual} du benz\'eno\"ide : chaque hexagone devient un sommet,
chaque ar\^ete partag\'ee devient une ar\^ete duale, et la taille de chaque
cycle devient une variable enti\`ere \`a domaine $\{5,6,7\}$. Cette
formulation permet :
\begin{itemize}
  \item d'exploiter la structure d'incidence du benz\'eno\"ide comme
        donn\'ee d'entr\'ee (le graphe dual n'est pas une inconnue) ;
  \item de filtrer agressivement les domaines par pr\'e-traitement
        avant la r\'esolution ;
  \item de tirer parti des outils \'eprouv\'es de la litt\'erature CSP
        (PyCSP3 pour la mod\'elisation, ACE pour la r\'esolution,
        \textsc{nauty}/VF2 pour le calcul d'automorphismes).
\end{itemize}

\subsection{Plan du chapitre}

La suite du chapitre est organis\'ee comme suit. La
section~\ref{sec:notations} fixe les notations et d\'efinitions
pr\'eliminaires. La section~\ref{sec:modelisation-graphe} introduit la
mod\'elisation par graphes : graphe de carbone, graphe dual, patterns
d'occupation, positions d'ar\^etes. La section~\ref{sec:propstructurelles}
\'etablit les propri\'et\'es structurelles utilis\'ees en pr\'e-traitement,
en particulier la caract\'erisation des hexagones \emph{gel\'es}. La
section~\ref{sec:csp} pr\'esente le mod\`ele CSP complet avec ses
contraintes C1 \`a~C5 et leurs preuves. La section~\ref{sec:symetrie}
traite de la rupture de sym\'etrie via le th\'eor\`eme de Whitney et les
contraintes lex-leader. La section~\ref{sec:reconstruction} d\'ecrit les
algorithmes de graphe (parcours BFS sur le dual, contraction/insertion)
utilis\'es pour la reconstruction tridimensionnelle des solutions
\'enum\'er\'ees par le CSP. La section~\ref{sec:pipeline} synth\'etise le
pipeline exp\'erimental et les complexit\'es. La
section~\ref{sec:conclusion} conclut.

% ════════════════════════════════════════════════════════════
\section{Notations et d\'efinitions pr\'eliminaires}
\label{sec:notations}
% ════════════════════════════════════════════════════════════

Nous fixons ici le vocabulaire utilis\'e dans tout le chapitre. Les notions
abord\'ees rel\`event de la \textbf{th\'eorie des graphes planaires} et
de la combinatoire \'el\'ementaire ; les lecteurs familiers peuvent
parcourir cette section rapidement.

\subsection{Graphes : rappels}

\begin{definition}[Graphe planaire]
Un graphe $G = (V, E)$ est \emph{planaire} s'il admet un plongement dans
le plan euclidien tel que les ar\^etes ne se croisent qu'en leurs
extr\'emit\'es. Un \emph{plongement planaire} d\'efinit un ensemble de
\emph{faces}, dont l'une (la face \emph{externe}) est non born\'ee.
\end{definition}

\begin{theoreme}[Formule d'Euler]
Pour tout graphe planaire connexe avec $V$ sommets, $E$ ar\^etes
et $F$ faces (face externe incluse) :
\[
  V - E + F = 2.
\]
\end{theoreme}

\begin{definition}[$k$-connexit\'e]
Un graphe $G$ est \emph{$k$-connexe} si $|V(G)| > k$ et si la suppression
de tout ensemble de $k-1$ sommets le laisse connexe. Un graphe est
\emph{2-connexe} (= connexe sans point d'articulation) si toute paire de
sommets distincts est connect\'ee par au moins deux chemins
sommet-disjoints.
\end{definition}

\subsection{Benz\'eno\"ides et substitution de cycles}

\begin{definition}[Benz\'eno\"ide]
Un \emph{benz\'eno\"ide} $\B$ est un sous-graphe fini, planaire, connexe
et 2-connexe du r\'eseau hexagonal r\'egulier, dont toutes les faces
internes sont des hexagones et dont chaque sommet a un degr\'e dans
$\{2, 3\}$.
\end{definition}

Nous notons :
\begin{itemize}
  \item $h$ : le nombre d'hexagones (= faces internes) de $\B$ ;
  \item $V$, $E$ : le nombre de sommets (atomes de carbone) et
        d'ar\^etes (liaisons C--C) de $\B$ ;
  \item $\ncinq$, $\nhex$, $\nsept$ : les nombres de cycles
        respectivement pentagonaux, hexagonaux, heptagonaux apr\`es
        substitution.
\end{itemize}

\begin{definition}[Substitution admissible]
Une \emph{substitution admissible} de $\B$ est une assignation
$\alpha : F(\B) \to \{5, 6, 7\}$ associant \`a chaque face hexagonale
interne une nouvelle taille de cycle, telle que la mol\'ecule r\'esultante
respecte la conservation des atomes de carbone et les contraintes
topologiques (d\'etaill\'ees au fil de ce chapitre).
\end{definition}

\subsection{Notations alg\'ebriques}

Le tableau~\ref{tab:notations} r\'esume les notations utilis\'ees.

\begin{table}[!ht]
\centering
\renewcommand{\arraystretch}{1.25}
\begin{tabular}{@{}lp{10.5cm}@{}}
\toprule
\textbf{Notation} & \textbf{Signification} \\
\midrule
$\B$ & Benz\'eno\"ide d'entr\'ee \\
$\GB = (V_\B, E_\B, F_\B)$ & Graphe de carbone de $\B$ : sommets, ar\^etes, faces \\
$\GD = (\VD, \ED)$ & Graphe dual de $\B$ \\
$h = |\VD|$ & Nombre d'hexagones du benz\'eno\"ide \\
$\xv$ & Variable taille du cycle $v \in \VD$, $\xv \in \{5,6,7\}$ \\
$\sigma(v) \in \{0,1\}^6$ & Pattern (motif d'occupation) des c\^ot\'es de $v$ \\
$\posIn{v}{e} \in \{0,\ldots,5\}$ & Position de l'ar\^ete $e$ dans le cycle $v$ \\
$\NN(v)$ & Voisinage de $v$ dans $\GD$ \\
$\degv$ & Degr\'e de $v$ dans $\GD$, $1 \leq \degv \leq 6$ \\
$b(v)$ & Nombre de blocs cycliques de z\'eros dans $\sigma(v)$ \\
$I(v)$ & Ensemble des sommets int\'erieurs libres de $v$ \\
$\Tn$ & Table des configurations de voisinage admissibles pour cycle de taille $n$ \\
$\Lambda(v)$ & Configuration locale de $v$ : $(\xv, x_{u_1}, \ldots, x_{u_k})$ \\
$\Aut(\GD)$ & Groupe des automorphismes de $\GD$ \\
$\dom(\xv)$ & Domaine courant de la variable $\xv$ \\
\bottomrule
\end{tabular}
\caption{Principales notations utilis\'ees dans ce chapitre.}
\label{tab:notations}
\end{table}

% ════════════════════════════════════════════════════════════
\section{Mod\'elisation par graphes}
\label{sec:modelisation-graphe}
% ════════════════════════════════════════════════════════════

Trois objets graph-th\'eoriques sont mis en jeu : le \emph{graphe de
carbone} $\GB$ (structure physique de la mol\'ecule), le \emph{graphe dual}
$\GD$ (objet central du mod\`ele CSP) et, plus tard, des graphes
auxiliaires utilis\'es pour la reconstruction 3D.

\subsection{Le graphe de carbone : entr\'ee}
\label{sec:graphe-carbone}

\subsubsection{D\'efinition formelle}

\begin{definition}[Graphe de carbone]
Soit $\B$ un benz\'eno\"ide. Le \emph{graphe de carbone} de $\B$ est le
triplet $\GB = (V_\B,\, E_\B,\, F_\B)$ o\`u :
\begin{itemize}
  \item $V_\B$ est l'ensemble des atomes de carbone ;
  \item $E_\B \subseteq \binom{V_\B}{2}$ est l'ensemble des liaisons C--C ;
  \item $F_\B$ est l'ensemble des hexagones (faces internes), chacun
        d\'ecrit par une suite ordonn\'ee de 6~sommets adjacents par paires.
\end{itemize}
\end{definition}

\begin{propriete}
Le graphe de carbone $\GB$ d'un benz\'eno\"ide est planaire, connexe,
2-connexe, biparti et de degr\'e maximum 3.
\end{propriete}

\subsubsection{Format d'entr\'ee Benzai}

Le format d'entr\'ee (\og Benzai \fg{}, inspir\'e de DIMACS) encode
$\GB$ sous forme textuelle. Le listing~\ref{lst:benzai} pr\'esente
l'exemple \texttt{first.graph}, un benz\'eno\"ide lin\'eaire \`a
3~hexagones (= 14 carbones).

\begin{lstlisting}[style=benzai,
  caption={Fichier Benzai d'un benz\'eno\"ide \`a 3 hexagones lin\'eaires.},
  label={lst:benzai}]
p DIMACS 14 16 3
e 0_0 1_1
e 0_0 -1_1
e 1_1 1_2
e -1_1 -1_2
e 1_2 0_3
...
h 0_0 1_1 1_2 0_3 -1_2 -1_1
h -1_2 0_3 0_4 -1_5 -2_4 -2_3
h -2_4 -1_5 -1_6 -2_7 -3_6 -3_5
\end{lstlisting}

\begin{itemize}
  \item Ligne \texttt{p DIMACS $V$ $E$ $F$} : ent\^ete avec
        $V$~sommets, $E$~ar\^etes, $F$~faces hexagonales ;
  \item Lignes \texttt{e $s_1$ $s_2$} : ar\^ete entre les sommets
        $s_1$ et $s_2$. Les identifiants $x\_y$ sont des coordonn\'ees
        dans le r\'eseau hexagonal ;
  \item Lignes \texttt{h $s_1$ $s_2$ $\ldots$ $s_6$} : hexagone d\'ecrit
        par ses 6~sommets dans l'\textbf{ordre cyclique}.
\end{itemize}

L'ordre cyclique des sommets dans chaque ligne \texttt{h} est crucial :
il fournit implicitement la \emph{num\'erotation} des c\^ot\'es de
l'hexagone (le c\^ot\'e $i$ est l'ar\^ete $(s_i, s_{i+1})$ modulo~6).

\begin{figure}[!ht]
\centering
\begin{tikzpicture}[scale=0.85, every node/.style={font=\small}]
  % Hexagone v0
  \foreach \i in {0,...,5} {
    \coordinate (A\i) at ({cos(60*\i+30)*1.2}, {sin(60*\i+30)*1.2});
  }
  \draw[thick,fill=grisHex] (A0) -- (A1) -- (A2) -- (A3) -- (A4) -- (A5) -- cycle;
  \node at (0,0) {\textbf{$v_0$}};
  \foreach \i in {0,...,5} {\fill (A\i) circle (2pt);}
  % Hexagone v1
  \foreach \i in {0,...,5} {
    \coordinate (B\i) at ({cos(60*\i+30)*1.2 + 2*1.2*cos(30)}, {sin(60*\i+30)*1.2});
  }
  \draw[thick,fill=grisHex] (B0) -- (B1) -- (B2) -- (B3) -- (B4) -- (B5) -- cycle;
  \node at ({2*1.2*cos(30)},0) {\textbf{$v_1$}};
  \foreach \i in {0,...,5} {\fill (B\i) circle (2pt);}
  % Hexagone v2
  \foreach \i in {0,...,5} {
    \coordinate (C\i) at ({cos(60*\i+30)*1.2 + 4*1.2*cos(30)}, {sin(60*\i+30)*1.2});
  }
  \draw[thick,fill=grisHex] (C0) -- (C1) -- (C2) -- (C3) -- (C4) -- (C5) -- cycle;
  \node at ({4*1.2*cos(30)},0) {\textbf{$v_2$}};
  \foreach \i in {0,...,5} {\fill (C\i) circle (2pt);}
\end{tikzpicture}
\caption{Le benz\'eno\"ide \texttt{first.graph} (3 hexagones lin\'eaires,
14 carbones). Chaque sommet noir est un atome de carbone ; chaque ar\^ete
est une liaison C--C. Les hexagones $v_0$, $v_1$, $v_2$ sont nos
\emph{faces} d'int\'er\^et.}
\label{fig:first-graph}
\end{figure}

\subsection{Le graphe dual}
\label{sec:dual}

\subsubsection{Construction et propri\'et\'es}

Le \textbf{graphe dual} est l'objet central de notre mod\'elisation. Il
encode la \emph{topologie de fusion} des hexagones et oublie les
coordonn\'ees individuelles des atomes.

\begin{definition}[Graphe dual]
\label{def:dual}
Soit $\GB = (V_\B, E_\B, F_\B)$ le graphe de carbone d'un benz\'eno\"ide
$\B$. Le \emph{graphe dual} $\GD = (\VD, \ED)$ est d\'efini par :
\begin{itemize}
  \item $\VD = \{0, 1, \ldots, h-1\}$, un sommet par hexagone de $F_\B$ ;
  \item $\{u, v\} \in \ED$ si et seulement si les hexagones $u$ et $v$
        partagent une ar\^ete dans $E_\B$.
\end{itemize}
\end{definition}

\begin{figure}[!ht]
\centering
\begin{tikzpicture}[scale=1, every node/.style={font=\small}]
  \node[circle,draw,thick,fill=bleuSub!15,minimum size=8mm] (v0) at (0,0) {$v_0$};
  \node[circle,draw,thick,fill=bleuSub!15,minimum size=8mm] (v1) at (2.5,0) {$v_1$};
  \node[circle,draw,thick,fill=bleuSub!15,minimum size=8mm] (v2) at (5,0) {$v_2$};
  \draw[thick] (v0) -- (v1);
  \draw[thick] (v1) -- (v2);
\end{tikzpicture}
\caption{Graphe dual $\GD$ du benz\'eno\"ide \texttt{first.graph} :
un chemin $v_0 - v_1 - v_2$. Le sommet $v_1$, de degr\'e~2, est le seul
hexagone partageant deux c\^ot\'es avec ses voisins.}
\label{fig:dual-first}
\end{figure}

\begin{propriete}[Propri\'et\'es du graphe dual]
\label{prop:dual}
Pour tout benz\'eno\"ide $\B$, son graphe dual $\GD$ satisfait :
\begin{enumerate}[label=(\roman*),leftmargin=*]
  \item $\GD$ est planaire ;
  \item le degr\'e de tout sommet $v \in \VD$ v\'erifie $1 \leq \degv \leq 6$ ;
  \item $\GD$ est connexe ssi $\B$ est connexe ;
  \item $\GD$ est un sous-graphe induit du graphe dual du \emph{pavage
        hexagonal infini} (lui-m\^eme isomorphe au \emph{r\'eseau
        triangulaire}).
\end{enumerate}
\end{propriete}

\begin{proof}
(i) Tout graphe est planaire ssi il est sous-graphe d'un graphe planaire ;
le pavage hexagonal infini est planaire, son graphe dual l'est aussi.
(ii)~Un hexagone du pavage admet au plus 6~voisins (un par c\^ot\'e), et
au moins~1 dans un benz\'eno\"ide connexe \`a $h \geq 2$ hexagones.
(iii)~\'Equivalence \'el\'ementaire entre connexit\'e de $\B$ et de
$\GD$. (iv)~D\'ecoule de la d\'efinition du dual du pavage hexagonal :
ses sommets sont les centres d'hexagones, et deux centres sont adjacents
ssi les hexagones partagent une ar\^ete. C'est pr\'ecis\'ement le
r\'eseau triangulaire.
\end{proof}

\subsubsection{Algorithme de construction}

La construction du dual proc\`ede par double boucle sur les paires
d'hexagones, avec test d'intersection des ensembles d'ar\^etes. La
complexit\'e est $\bigO(h^2 \cdot c)$ o\`u $c$ est le co\^ut de
l'intersection (constant si l'on indexe les ar\^etes par hash).
Pour $h \leq 20$ (cas de nos exp\'eriences), c'est instantan\'e.

\begin{algorithm}[!ht]
\caption{Construction du graphe dual $\GD$ depuis $\GB$.}
\label{alg:dual}
\begin{algorithmic}[1]
\Require Graphe de carbone $\GB = (V_\B, E_\B, F_\B)$, hexagones num\'erot\'es
$0, \ldots, h-1$ avec leurs listes d'ar\^etes $E(v)$.
\Ensure Graphe dual $\GD = (\VD, \ED)$ et positions $\posIn{\cdot}{\cdot}$.
\State $\VD \gets \{0, \ldots, h-1\}$ ;\quad $\ED \gets \emptyset$
\For{$u = 0$ \textbf{to} $h-1$}
  \For{$v = u+1$ \textbf{to} $h-1$}
    \State $S \gets E(u) \cap E(v)$ \Comment{ar\^etes partag\'ees}
    \If{$S \neq \emptyset$}
      \State $e_\text{shared} \gets$ choisir un \'el\'ement de $S$
      \State $\ED \gets \ED \cup \{(u, v)\}$
      \State $\posIn{u}{(u,v)} \gets$ position de $e_\text{shared}$ dans la s\'equence de $u$
      \State $\posIn{v}{(u,v)} \gets$ position de $e_\text{shared}$ dans la s\'equence de $v$
    \EndIf
  \EndFor
\EndFor
\State \Return $\GD, \posIn{\cdot}{\cdot}$
\end{algorithmic}
\end{algorithm}

\subsection{Patterns d'occupation et positions d'ar\^etes}
\label{sec:patterns}

Le graphe dual capture la topologie globale, mais la d\'ecision \og ce
c\^ot\'e de l'hexagone est-il partag\'e ? \fg{} requiert une information
\emph{locale} suppl\'ementaire : le \emph{pattern d'occupation}.

\begin{definition}[Pattern d'occupation]
\label{def:pattern}
Pour chaque sommet $v \in \VD$, le \emph{pattern d'occupation}
$\sigma(v) = (\sigma_0, \sigma_1, \ldots, \sigma_5) \in \{0, 1\}^6$
est d\'efini par
\[
  \sigma_i = \begin{cases}
    1 & \text{si le c\^ot\'e $i$ de $v$ est partag\'e avec un voisin dans $\GD$,}\\
    0 & \text{sinon.}
  \end{cases}
\]
\end{definition}

Le pattern $\sigma(v)$ est consid\'er\'e \emph{cycliquement} : les positions
$0, 1, 2, 3, 4, 5$ sont arrang\'ees autour de l'hexagone dans le sens
horaire, et la rotation $\sigma_i \mapsto \sigma_{i+1 \bmod 6}$ correspond
\`a une rotation g\'eom\'etrique de l'hexagone.

\begin{definition}[Position d'une ar\^ete]
Pour chaque ar\^ete $e = (u, v) \in \ED$, les fonctions
$\posIn{u}{e} \in \{0, \ldots, 5\}$ et $\posIn{v}{e} \in \{0, \ldots, 5\}$
indiquent les c\^ot\'es respectifs des hexagones $u$ et $v$ qui forment
l'ar\^ete physique partag\'ee.
\end{definition}

\begin{exemple}
Pour le benz\'eno\"ide \texttt{first.graph} :
\begin{center}
\renewcommand{\arraystretch}{1.2}
\begin{tabular}{ccl}
\toprule
Hexagone & $\degv$ & Pattern $\sigma(v)$ \\
\midrule
$v_0$ & 1 & $(0, 0, 0, 1, 0, 0)$ \\
$v_1$ & 2 & $(1, 0, 0, 1, 0, 0)$ \\
$v_2$ & 1 & $(1, 0, 0, 0, 0, 0)$ \\
\bottomrule
\end{tabular}
\end{center}
$v_0$ et $v_2$ sont des \og hexagones de bord \fg{} (un seul voisin).
$v_1$ est le seul hexagone central avec deux c\^ot\'es partag\'es.
\end{exemple}

\begin{propriete}[Consistance aux jonctions]
\label{prop:consistance}
Pour toute ar\^ete $e = (u, v) \in \ED$, les positions
$\posIn{u}{e}$ et $\posIn{v}{e}$ d\'eterminent une ar\^ete physique unique de
$\GB$. R\'eciproquement, deux hexagones $u, v$ adjacents dans $\GD$
partagent \emph{exactement une} ar\^ete dans $\GB$.
\end{propriete}

Cette propri\'et\'e est cruciale : elle implique que le \emph{nombre
d'ar\^etes du dual} \'egale le \emph{nombre d'ar\^etes partag\'ees}
dans le benz\'eno\"ide, et nous l'utilisons dans la preuve de la
contrainte~C1.

\subsection{Blocs cycliques et structure du voisinage}

Le pattern $\sigma(v)$ se d\'ecompose naturellement en \emph{blocs} de
positions cons\'ecutives valant~$0$ (c\^ot\'es libres) ou~$1$ (c\^ot\'es
partag\'es). Ces blocs jouent un r\^ole central dans la caract\'erisation
des hexagones immobiles (section~\ref{sec:propstructurelles}).

\begin{definition}[Blocs de z\'eros]
\label{def:blocs}
Soit $\sigma(v) \in \{0,1\}^6$. Un \emph{bloc de z\'eros} de $\sigma(v)$
est une sous-suite cyclique maximale de positions cons\'ecutives
valant~$0$. On note $b(v)$ le nombre de tels blocs.
\end{definition}

\begin{exemple}
\begin{itemize}
  \item $\sigma = (1, 1, 0, 0, 0, 0)$ : un seul bloc de 4~z\'eros,
        $b = 1$.
  \item $\sigma = (1, 0, 0, 1, 0, 0)$ : deux blocs de 2~z\'eros chacun
        ($\sigma_1\sigma_2$ et $\sigma_4\sigma_5$), $b = 2$.
  \item $\sigma = (1, 0, 1, 0, 1, 0)$ : trois blocs de 1~z\'ero,
        $b = 3$.
  \item $\sigma = (1, 1, 1, 1, 1, 1)$ : aucun z\'ero, $b = 0$.
\end{itemize}
\end{exemple}

L'algorithme~\ref{alg:blocs} calcule $b(v)$ en temps lin\'eaire en partant
d'une position partag\'ee (afin de ne pas couper un bloc cyclique).

\begin{algorithm}[!ht]
\caption{Calcul du nombre de blocs de z\'eros $b(v)$.}
\label{alg:blocs}
\begin{algorithmic}[1]
\Require Pattern $\sigma(v) \in \{0,1\}^6$ tel que $\degv \geq 1$.
\Ensure $b(v) \in \{0, 1, 2, 3\}$.
\State $start \gets$ premier indice $i$ tel que $\sigma_i = 1$
\State $blocs \gets 0$ ;\quad $in\_zero \gets \mathrm{faux}$
\For{$\text{offset} = 0$ \textbf{to} $5$}
  \State $idx \gets (start + \text{offset}) \bmod 6$
  \If{$\sigma_{idx} = 0$ \textbf{et} $\lnot in\_zero$}
    \State $blocs \gets blocs + 1$ ;\quad $in\_zero \gets \mathrm{vrai}$
  \ElsIf{$\sigma_{idx} = 1$}
    \State $in\_zero \gets \mathrm{faux}$
  \EndIf
\EndFor
\State \Return $blocs$
\end{algorithmic}
\end{algorithm}

% ════════════════════════════════════════════════════════════
\section{Propri\'et\'es structurelles et pr\'e-traitement}
\label{sec:propstructurelles}
% ════════════════════════════════════════════════════════════

Avant de poser le mod\`ele CSP complet, nous \'etablissons trois familles
de propri\'et\'es \og locales \fg{} qui permettent de \textbf{r\'eduire
les domaines des variables avant la r\'esolution}. Ces propri\'et\'es
sont calcul\'ees en \emph{temps constant} par sommet, \`a partir du seul
pattern $\sigma(v)$. Elles incarnent l'avantage central de notre
mod\'elisation : exploiter la g\'eom\'etrie de l'entr\'ee pour all\'eger
le travail du solveur.

\subsection{Hexagones int\'erieurs : la r\`egle $\degv = 6$}

\begin{propriete}[Hexagone enti\`erement entour\'e]
\label{prop:degre6}
Si $\degv = 6$, alors $\sigma(v) = (1,1,1,1,1,1)$ et le cycle $v$ ne peut
\^etre substitu\'e : on doit avoir $\xv = 6$.
\end{propriete}

\begin{proof}
$\degv = 6$ signifie que les six c\^ot\'es de $v$ sont partag\'es. Toute
substitution par un pentagone (contraction d'un sommet) ou un heptagone
(insertion d'un sommet) toucherait n\'ecessairement \`a un c\^ot\'e
partag\'e, donc cassait la topologie d'un cycle voisin. La seule
op\'eration compatible est l'identit\'e (cycle reste hexagonal).
\end{proof}

\subsection{Hexagones \`a blocs s\'epar\'es : la r\`egle $b(v) \geq 2$}

\begin{propriete}[Hexagone fig\'e par $b(v) \geq 2$]
\label{prop:b2}
Si $b(v) \geq 2$, alors le cycle $v$ ne peut \^etre substitu\'e : on doit
avoir $\xv = 6$.
\end{propriete}

\begin{proof}
Supposons $b(v) \geq 2$. Le bord de $v$ se d\'ecompose en au moins deux
blocs disjoints d'ar\^etes libres, s\'epar\'es par des blocs d'ar\^etes
partag\'ees. La cha\^ine cyclique des ar\^etes alterne donc :
\[
  [\text{partag\'ees}_1] - [\text{libres}_1] - [\text{partag\'ees}_2]
  - [\text{libres}_2] - \cdots
\]
Substituer $v$ par un pentagone n\'ecessite de \emph{contracter}
exactement un sommet (perte d'une ar\^ete libre) ; cela r\'eduirait la
longueur d'\emph{un seul} bloc libre tout en pr\'eservant la position
relative des autres ar\^etes partag\'ees. Or les blocs partag\'es sont
\emph{verrouill\'es} aux positions impos\'ees par les cycles voisins :
les distances cycliques entre eux sont fix\'ees. R\'eduire la longueur
totale du bord (de 6~\`a~5) sans changer ces distances est arithm\'etiquement
impossible.

Sym\'etriquement, substituer $v$ par un heptagone ajouterait une ar\^ete
libre dans un seul bloc, augmentant la longueur d'un seul intervalle mais
laissant les autres inchang\'es. Or la nouvelle somme des longueurs
(=~7) doit \^etre coh\'erente avec un polygone r\'egulier, ce qui force
toutes les distances cycliques \`a \^etre celles d'un heptagone. Impossible
si $b(v) \geq 2$.

Dans le cas $b(v) = 1$, en revanche, toutes les ar\^etes libres forment
un \emph{unique} intervalle cyclique : on peut l'agrandir (heptagone) ou
le r\'eduire (pentagone) sans violer la position des c\^ot\'es partag\'es.
\end{proof}

\begin{remarque}
Le crit\`ere $\degv = 6 \lor b(v) \geq 2$ est purement structurel : il se
d\'eduit du seul pattern $\sigma(v)$ et peut \^etre \'evalu\'e en
$\bigO(1)$ par sommet. C'est un cas typique de \emph{propagation
unaire} en CSP : la variable est r\'eduite \`a un singleton avant m\^eme
le lancement du solveur.
\end{remarque}

\subsection{R\'ealisabilit\'e topologique : sommets int\'erieurs et c\^ot\'es libres}
\label{sec:c2bis-prep}

Les propri\'et\'es ci-dessus ne suffisent pas. Sur des benz\'eno\"ides
\`a partir de $h \geq 6$ hexagones, nous avons observ\'e \emph{empiriquement}
que de nombreuses assignations satisfaisant C1+C2+C3+C4 produisaient des
\'echecs lors de la phase de reconstruction tridimensionnelle (\og
solutions g\'eom\'etriquement infaisables \fg{}, jusqu'\`a 75~\% sur
certaines mol\'ecules de $h_9$). Cela r\'ev\`ele une condition n\'ecessaire
\emph{suppl\'ementaire}, encod\'ee par les notions suivantes.

\begin{definition}[Sommet int\'erieur libre]
\label{def:interieur}
Soit $v \in \VD$ et $\sigma(v) = (\sigma_0, \ldots, \sigma_5)$. Un
\emph{sommet int\'erieur libre} de $v$ est un indice $i \in \{0,\ldots,5\}$
de l'hexagone tel que les deux ar\^etes incidentes (positions $i-1$
et~$i$, modulo~6) soient toutes deux libres :
\[
  \sigma_{(i-1) \bmod 6} = 0 \;\text{ et }\; \sigma_i = 0.
\]
On note $I(v) \subseteq \{0, \ldots, 5\}$ l'ensemble de ces sommets.
\end{definition}

\begin{definition}[C\^ot\'e libre]
Un \emph{c\^ot\'e libre} de $v$ est un indice $i$ tel que $\sigma_i = 0$.
L'ensemble des c\^ot\'es libres est $F_v = \{i : \sigma_i = 0\}$.
\end{definition}

\begin{propriete}[Condition n\'ecessaire de substitution]
\label{prop:c2bis}
Pour tout sommet $v \in \VD$ non gel\'e par les
Propri\'et\'es~\ref{prop:degre6}--\ref{prop:b2} :
\begin{enumerate}[label=(\alph*)]
  \item $\xv = 5$ est topologiquement r\'ealisable
        \emph{seulement si} $I(v) \neq \emptyset$ ;
  \item $\xv = 7$ est topologiquement r\'ealisable
        \emph{seulement si} $F_v \neq \emptyset$.
\end{enumerate}
\end{propriete}

\begin{proof}
\textbf{(a)} La substitution $\xv = 6 \to \xv = 5$ se r\'ealise
g\'eom\'etriquement par \emph{contraction} d'un sommet de l'hexagone :
on identifie deux sommets adjacents, ce qui revient \`a fusionner
deux ar\^etes incidentes en une seule. Pour que cette fusion ne casse
aucun cycle voisin, il faut que les \emph{deux} ar\^etes \`a fusionner
soient libres (i.e. non partag\'ees). Cela revient \`a exiger l'existence
d'un sommet $i \in I(v)$.

R\'eciproquement, si $I(v) \neq \emptyset$, soit $i \in I(v)$ ; la
contraction de ce sommet d\'efinit un pentagone bien form\'e dont les
c\^ot\'es partag\'es de $v$ restent inchang\'es. La phase de placement
g\'eom\'etrique (section~\ref{sec:reconstruction}) construit alors une
g\'eom\'etrie valide.

\textbf{(b)} La substitution $\xv = 6 \to \xv = 7$ se r\'ealise par
\emph{insertion} d'un nouveau sommet \emph{au milieu} d'une ar\^ete.
Cette insertion ne touche pas aux cycles voisins ssi l'ar\^ete choisie
est libre. La condition $F_v \neq \emptyset$ est donc n\'ecessaire et
suffisante.
\end{proof}

\begin{exemple}[Pattern $\sigma = (1,1,1,1,1,0)$]
$\degv = 5$, $b(v) = 1$, donc non gel\'e par C2. Mais :
\begin{itemize}
  \item L'unique $0$ est en position~5, et ses deux voisins cycliques
        (positions 4 et 0) valent~1. Donc $I(v) = \emptyset$.
  \item $F_v = \{5\}$, donc $F_v \neq \emptyset$.
\end{itemize}
Conclusion : $\dom(\xv) = \{6, 7\}$ ; le pentagone est interdit par
notre Propri\'et\'e~\ref{prop:c2bis}.
\end{exemple}

\subsection{Calcul des domaines}

L'algorithme~\ref{alg:domaines} combine les trois r\`egles pour produire
le \emph{domaine initial} de chaque variable. Sa complexit\'e est
$\bigO(h)$.

\begin{algorithm}[!ht]
\caption{Calcul des domaines initiaux $\dom(\xv)$.}
\label{alg:domaines}
\begin{algorithmic}[1]
\Require Graphe dual $\GD$, patterns $\sigma(v)$, indicateur \texttt{freeze\_b2}.
\Ensure Domaines $\dom(\xv)$ pour tout $v \in \VD$.
\For{$v \in \VD$}
  \If{$\degv = 6$} \Comment{R\`egle 1 : enti\`erement entour\'e}
    \State $\dom(\xv) \gets \{6\}$
  \ElsIf{$\texttt{freeze\_b2} \land b(v) \geq 2$} \Comment{R\`egle 2}
    \State $\dom(\xv) \gets \{6\}$
  \Else \Comment{R\`egle 3 : C2bis}
    \State $\dom(\xv) \gets \{6\}$
    \If{$I(v) \neq \emptyset$}
      \State $\dom(\xv) \gets \dom(\xv) \cup \{5\}$
    \EndIf
    \If{$F_v \neq \emptyset$}
      \State $\dom(\xv) \gets \dom(\xv) \cup \{7\}$
    \EndIf
  \EndIf
\EndFor
\State \Return $\{\dom(\xv) : v \in \VD\}$
\end{algorithmic}
\end{algorithm}

\subsection{Tables de voisinage : structure et filtrage}

Les tables $\mathcal{T}(5)$, $\mathcal{T}(6)$, $\mathcal{T}(7)$
constituent une donn\'ee \emph{externe} au CSP, pr\'e-calcul\'ee par
validation g\'eom\'etrique de toutes les configurations \og cycle central
+ voisins \fg{} possibles. Leur construction est document\'ee dans le
pipeline g\'en\'eral (annexe du rapport principal) ; nous en r\'esumons
ici les caract\'eristiques.

\begin{definition}[Table de voisinage]
Pour $n \in \{5, 6, 7\}$, la table $\mathcal{T}(n)$ est l'ensemble des
tuples $(n, s_1, s_2, \ldots, s_n) \in \{5,6,7\}^{n+1}$ tels qu'il existe
une g\'eom\'etrie 3D \emph{planaire} d'un cycle central de taille $n$
entour\'e des cycles de tailles $s_1, \ldots, s_n$ (modulo rotation et
r\'eflexion du voisinage).
\end{definition}

Les tailles obtenues exp\'erimentalement, apr\`es validation
g\'eom\'etrique par xTB et test de planarit\'e, sont r\'ecapitul\'ees dans
le tableau~\ref{tab:tables-tailles}.

\begin{table}[!ht]
\centering
\renewcommand{\arraystretch}{1.2}
\begin{tabular}{lrr}
\toprule
Cycle central $n$ & $|\mathcal{T}(n)|$ (protocole xTB \texttt{--opt}) & $|\mathcal{T}(n)|$ (protocole MD + opt) \\
\midrule
5 &  92 &  63 \\
6 & 412 & 280 \\
7 & 699 & 367 \\
\midrule
\textbf{Total} & \textbf{1203} & \textbf{710} \\
\bottomrule
\end{tabular}
\caption{Tailles des tables de voisinage selon le protocole de validation.}
\label{tab:tables-tailles}
\end{table}

\subsubsection{Filtrage par pattern}

Pour chaque sommet $v$ du benz\'eno\"ide d'entr\'ee, la table \og
brute \fg{} $\mathcal{T}(n)$ contient beaucoup d'entr\'ees \emph{non
compatibles} avec la topologie locale de $v$ (i.e. avec son pattern
$\sigma(v)$). Nous restreignons donc la table \`a une \emph{table
locale} $\mathcal{T}_v(n)$.

\begin{definition}[Table restreinte]
$\mathcal{T}_v(n) \subseteq \mathcal{T}(n)$ est l'ensemble des tuples
$(n, s_1, \ldots, s_n) \in \mathcal{T}(n)$ tels que :
\begin{enumerate}[label=(\roman*)]
  \item le nombre de positions $i$ avec $s_i \neq 0$ vaut $\degv$ ;
  \item les positions $i$ avec $s_i \neq 0$ forment un bloc consécutif
        cyclique (puisque $b(v) = 1$ pour les sommets non gel\'es) ;
  \item les valeurs $s_i$ aux positions partag\'ees sont compatibles avec
        les domaines courants $\dom(x_{u_j})$ des voisins $u_j$ de $v$.
\end{enumerate}
\end{definition}

\subsubsection{Arc-consistance et propagation}

Les tables restreintes alimentent un calcul d'\emph{arc-consistance binaire}
classique : pour chaque ar\^ete $(u, v) \in \ED$, on \'elimine de
$\mathcal{T}_v$ toute entr\'ee dont la valeur \`a la position de $u$ n'est
pas dans $\dom(x_u)$, et r\'eciproquement. Ce processus se propage jusqu'\`a
point fixe.

\begin{propriete}[Complexit\'e du filtrage AC]
La proc\'edure d'arc-consistance binaire sur le couple
$(\GD, \mathcal{T})$ termine en temps
$\bigO(|\ED| \cdot |\mathcal{T}|_{\max})$
o\`u $|\mathcal{T}|_{\max} = \max_n |\mathcal{T}(n)|$.
\end{propriete}

% ════════════════════════════════════════════════════════════
\section{Le mod\`ele CSP}
\label{sec:csp}
% ════════════════════════════════════════════════════════════

\subsection{Structure g\'en\'erale du mod\`ele}

Le mod\`ele CSP est d\'efini par un triplet $(X, D, C)$ o\`u $X$ est
l'ensemble des variables, $D$ leurs domaines, et $C$ l'ensemble des
contraintes. Pour notre probl\`eme :
\[
  X = \{\xv : v \in \VD\}, \qquad
  D = \{\dom(\xv) \subseteq \{5, 6, 7\} : v \in \VD\}, \qquad
  C = \{C1, C2, C2\text{bis}, C3, C4, C5\}.
\]

\subsection{Variables et domaines}

\begin{definition}[Variables]
Pour chaque sommet $v \in \VD$, nous d\'efinissons une variable
$\xv$ \`a domaine $\dom(\xv) \subseteq \{5, 6, 7\}$, repr\'esentant la
\emph{taille} du cycle au sommet $v$ apr\`es substitution.
\end{definition}

L'ensemble des variables est $\mathbf{X} = (\xv)_{v \in \VD}$ avec
$|\mathbf{X}| = h$.

\subsection{Contrainte C1 : conservation du carbone}

\begin{contrainte}[Conservation]
\label{c:carbone}
\[
  \sum_{v \in \VD} \xv \;=\; 6\,h.
\]
\end{contrainte}

Cette contrainte exprime la conservation du nombre total d'atomes de
carbone au cours d'une substitution. Sa justification rigoureuse passe
par la formule d'Euler.

\begin{propriete}[Formule d'Euler appliqu\'ee \`a $\GB$]
\label{prop:euler}
Pour un syst\`eme fusionn\'e de $h$~cycles formant un graphe planaire
connexe avec $V$ sommets et $E$ ar\^etes :
\[
  V \;=\; \sum_{v \in \VD} \xv \;-\; |\ED| \;-\; h \;+\; 1,
\]
o\`u $|\ED|$ est le nombre d'ar\^etes du graphe dual.
\end{propriete}

\begin{proof}
La formule d'Euler pour un graphe planaire connexe avec $V$ sommets,
$E$ ar\^etes et $F$ faces (face externe incluse) s'\'ecrit
$V - E + F = 2$. Nous avons $F = h + 1$ (les $h$ cycles internes plus
la face externe), donc
\[
  V = E - h + 1.
\]
Il reste \`a exprimer $E$ en fonction des tailles des cycles. Chaque
cycle $v$ poss\`ede $\xv$ ar\^etes ; mais une ar\^ete \emph{partag\'ee}
entre deux cycles est compt\'ee deux fois dans la somme
$\sum_{v} \xv$. Par la Propri\'et\'e~\ref{prop:consistance}, le nombre
d'ar\^etes partag\'ees \'egale $|\ED|$. Une ar\^ete \emph{de bord} (non
partag\'ee) est compt\'ee une seule fois. Donc :
\[
  \sum_{v \in \VD} \xv \;=\; E + |\ED|,
  \qquad\text{soit}\qquad
  E \;=\; \sum_{v \in \VD} \xv - |\ED|.
\]
Substituant dans $V = E - h + 1$ : $V = \sum \xv - |\ED| - h + 1$.
\end{proof}

\begin{corollaire}[\'Equivalence C1 $\Leftrightarrow$ $\ncinq = \nsept$]
La contrainte C1 est \'equivalente \`a $\ncinq = \nsept$.
\end{corollaire}

\begin{proof}
Posant $\ncinq + \nhex + \nsept = h$ et
$\sum \xv = 5\ncinq + 6\nhex + 7\nsept$, l'\'egalit\'e
$\sum \xv = 6h$ \'equivaut \`a
$5\ncinq + 6\nhex + 7\nsept = 6(\ncinq + \nhex + \nsept)$,
soit $-\ncinq + \nsept = 0$.
\end{proof}

\begin{remarque}[Conservation des atomes]
Puisque $|\ED|$ et $h$ sont des \emph{constantes} du benz\'eno\"ide
d'entr\'ee, la Propri\'et\'e~\ref{prop:euler} montre que
$V_{\text{substituant}} = V_{\text{original}}$ ssi
$\sum \xv = 6h$. C'est la formulation chimique de C1 : conservation
du nombre de carbones lors de la substitution.
\end{remarque}

\subsection{Contrainte C2 : hexagones gel\'es}

\begin{contrainte}[Gel des hexagones immobiles]
\label{c:gele}
\[
  \forall v \in \VD,\quad
  [\, \degv = 6 \;\lor\; b(v) \geq 2 \,]
  \;\Longrightarrow\; \xv = 6.
\]
\end{contrainte}

La justification compl\`ete est fournie par les
Propri\'et\'es~\ref{prop:degre6} et~\ref{prop:b2} de la
section~\ref{sec:propstructurelles}. Concr\`etement, l'impl\'ementation
n'ajoute pas une contrainte au sens classique : elle \emph{r\'eduit le
domaine} de la variable \`a $\{6\}$ avant le lancement du solveur.

\subsection{Contrainte C2bis : filtrage topologique du domaine}

\begin{contrainte}[Filtrage topologique]
\label{c:c2bis}
Pour tout $v \in \VD$ non gel\'e par C2 :
\[
  5 \in \dom(\xv) \;\Longleftrightarrow\; I(v) \neq \emptyset,
  \qquad
  7 \in \dom(\xv) \;\Longleftrightarrow\; F_v \neq \emptyset.
\]
\end{contrainte}

Cette contrainte, justifi\'ee par la Propri\'et\'e~\ref{prop:c2bis},
est appliqu\'ee comme C2 par r\'eduction de domaine au pr\'e-traitement
en $\bigO(h)$.

\subsection{Contrainte C3 : voisinage admissible}

\begin{definition}[Configuration locale]
La \emph{configuration locale} de $v$ est le tuple
\[
  \Lambda(v) \;=\; \bigl(\xv,\, x_{u_1}, x_{u_2}, \ldots, x_{u_k}\bigr),
\]
o\`u $(u_1, \ldots, u_k)$ est la suite ordonn\'ee des voisins de $v$
dans $\GD$ selon les positions $\posIn{v}{e}$ croissantes.
\end{definition}

\begin{contrainte}[Voisinage admissible]
\label{c:voisinage}
\[
  \forall v \in \VD,\quad \Lambda(v) \in \mathcal{T}(\xv).
\]
\end{contrainte}

C3 est une \textbf{contrainte extensionnelle} (au sens CSP) : elle
\'enum\`ere explicitement les tuples valides. Sa propagation est
efficace dans les solveurs CSP modernes (algorithme GAC-4, etc.).

\begin{remarque}[Sommets de degr\'e 1]
La table $\mathcal{T}(n)$ est construite avec un voisinage minimum de
2~cycles. Pour les sommets $v \in \VD$ de degr\'e 1, la contrainte C3
n'est pas appliqu\'ee : un cycle isol\'e d'un seul c\^ot\'e est toujours
g\'eom\'etriquement admissible.
\end{remarque}

\subsection{Contrainte C4 : brisure de sym\'etrie}

\begin{contrainte}[Lex-leader]
\label{c:lex}
Pour chaque g\'en\'erateur $\pi$ du groupe d'automorphismes
$\Aut(\GD)$ :
\[
  (\xv)_{v \in \VD}
  \;\leq_{\mathrm{lex}}\;
  (x_{\pi(v)})_{v \in \VD}.
\]
\end{contrainte}

La justification compl\`ete de cette contrainte (Th\'eor\`eme de
Whitney + caract\'erisation des redondances) est d\'evelopp\'ee
\`a la section~\ref{sec:symetrie}.

\subsection{Contrainte C5 : adjacence 5--7 (optionnelle)}

Cette contrainte est \emph{exp\'erimentale}, motiv\'ee par l'observation
chimique selon laquelle les pentagones et heptagones tendent \`a appara\^itre
en paires adjacentes (\emph{motifs azul\'eniques}).

\begin{contrainte}[Adjacence 5--7]
\label{c:adj57}
\[
  \forall v \in \VD,\quad
  \xv = 5 \;\Longrightarrow\; \exists u \in \NN(v),\; x_u = 7
\]
\[
  \forall v \in \VD,\quad
  \xv = 7 \;\Longrightarrow\; \exists u \in \NN(v),\; x_u = 5.
\]
\end{contrainte}

C5 \emph{n'interdit pas} les paires 5--5 ou 7--7 ; elle exige seulement
la pr\'esence d'au moins un voisin compl\'ementaire. Elle est
d\'esactiv\'ee par d\'efaut.

\subsection{R\'ecapitulatif du mod\`ele}

\begin{table}[!ht]
\centering
\renewcommand{\arraystretch}{1.4}
\begin{tabular}{@{}cll@{}}
\toprule
\textbf{Code} & \textbf{Nom} & \textbf{Formulation} \\
\midrule
C1 & Conservation du carbone &
  $\sum_{v \in \VD} \xv = 6h \;\Leftrightarrow\; \ncinq = \nsept$ \\
C2 & Hexagones gel\'es &
  $[\degv = 6 \lor b(v) \geq 2] \Rightarrow \xv = 6$ \\
C2bis & Filtrage topologique &
  $5 \in \dom(\xv) \Leftrightarrow I(v) \neq \emptyset$ \\
   & du domaine &
  $7 \in \dom(\xv) \Leftrightarrow F_v \neq \emptyset$ \\
C3 & Voisinage admissible &
  $\Lambda(v) \in \mathcal{T}(\xv)$ \\
C4 & Brisure de sym\'etrie &
  $\mathbf{x} \leq_{\mathrm{lex}} \pi(\mathbf{x})$ pour chaque g\'en\'erateur $\pi$ \\
C5 & Adjacence 5--7 &
  Implications d'existence d'un voisin de type compl\'ementaire \\
   & (optionnelle) & \\
\bottomrule
\end{tabular}
\caption{R\'ecapitulatif des contraintes du mod\`ele CSP.}
\label{tab:contraintes}
\end{table}

\subsection{Complexit\'e du mod\`ele et r\'esolution}

\subsubsection{Taille du probl\`eme}

\begin{itemize}
  \item Variables : $|\mathbf{X}| = h$, typiquement entre 2 et 20 dans
        nos exp\'eriences.
  \item Domaines : $|\dom(\xv)| \leq 3$ ; apr\`es pr\'e-traitement,
        souvent $\leq 2$ ou r\'eduit \`a un singleton.
  \item Espace de recherche brut : $\leq 3^h \approx 3{,}5 \cdot 10^9$
        pour $h = 20$.
\end{itemize}

\subsubsection{Choix du solveur}

Nous utilisons \textbf{PyCSP3} (biblioth\`eque Python g\'en\'erant des
fichiers XCSP3) coupl\'e \`a \textbf{ACE} (solveur Java tournant en
\emph{boite noire}). Le choix est motiv\'e par :
\begin{itemize}
  \item l'int\'egration native Python (pr\'e/post-traitement dans le m\^eme
        \'ecosyst\`eme) ;
  \item le format XCSP3 standardis\'e (interchangeabilit\'e des solveurs) ;
  \item le support des contraintes extensionnelles (C3) et lex (C4) en
        natif.
\end{itemize}

\subsubsection{Strat\'egie de propagation recommand\'ee}

L'ordre suivant minimise la taille de l'espace de recherche transmis au
solveur :
\[
  \text{C1} \to \text{C2} \to \text{C2bis} \to \text{AC}(\text{C3}) \to \text{C4} \to (\text{C5}).
\]
C1 globalise les contraintes de comptage ; C2 et C2bis r\'eduisent les
domaines \emph{une fois} en pr\'e-traitement ; l'arc-consistance sur C3
propage les r\'eductions de voisinage ; C4 \'elimine les doublons par
sym\'etrie ; C5 (optionnelle) restreint au motif azul\'enique.

% ════════════════════════════════════════════════════════════
\section{Sym\'etrie : automorphismes et brisure}
\label{sec:symetrie}
% ════════════════════════════════════════════════════════════

\subsection{Probl\`eme des solutions redondantes}

Le graphe dual $\GD$ admet en g\'en\'eral des \textbf{automorphismes}
non triviaux : permutations de ses sommets qui pr\'eservent l'adjacence.
Par exemple, un benz\'eno\"ide lin\'eaire \`a $h = 3$ hexagones admet un
unique automorphisme non identit\'e (la r\'eflexion qui \'echange les
deux hexagones de bord).

Si $\pi \in \Aut(\GD)$ et $\alpha : \VD \to \{5,6,7\}$ est une solution
du CSP, alors la solution $\beta = \pi \cdot \alpha$ d\'efinie par
$\beta(v) = \alpha(\pi^{-1}(v))$ d\'ecrit la \emph{m\^eme structure
mol\'eculaire} (\`a rotation et r\'eflexion pr\`es). Ces solutions ne
doivent \^etre compt\'ees qu'une fois.

La question \emph{th\'eorique} est : \og les automorphismes de $\GD$
capturent-ils toutes les redondances, et \emph{seulement} elles ? \fg{}
La r\'eponse est positive, et nous l'\'etablissons rigoureusement
ci-dessous.

\subsection{Groupe d'automorphismes}

\begin{definition}[Automorphisme du dual]
Un \emph{automorphisme} de $\GD$ est une bijection $\pi : \VD \to \VD$
telle que pour tout couple $(u, v)$ :
\[
  \{u, v\} \in \ED \;\Longleftrightarrow\; \{\pi(u), \pi(v)\} \in \ED.
\]
L'ensemble des automorphismes, muni de la composition, forme le
\emph{groupe d'automorphismes} $\Aut(\GD)$.
\end{definition}

\begin{definition}[Action sur les solutions]
$\Aut(\GD)$ agit sur l'ensemble des solutions par
$(\pi \cdot \alpha)(v) = \alpha(\pi^{-1}(v))$. Les \emph{orbites} de
cette action sont des classes d'\'equivalence de solutions
\og structurellement \'equivalentes \fg{}.
\end{definition}

\subsubsection{Calcul des automorphismes}

Calculer $\Aut(\GD)$ revient \`a r\'esoudre le probl\`eme d'\emph{
isomorphisme de graphes} sur $(\GD, \GD)$. C'est un probl\`eme classique
de complexit\'e \emph{quasi-polynomiale} (Babai, 2016) ; en pratique des
algorithmes tr\`es efficaces existent :
\begin{itemize}
  \item \textbf{VF2} (Cordella et al., 2004), impl\'ement\'e dans
        NetworkX : exploration par retour arri\`ere avec coupes par
        invariants (degr\'e, voisinage). Suffisant pour $h \leq 30$.
  \item \textbf{\textsc{nauty}/\textsc{bliss}} : exhaustif et tr\`es
        rapide en pratique (algorithme de partition par raffinement).
\end{itemize}

Pour les benz\'eno\"ides exp\'erimentaux ($h \leq 20$), VF2 est
suffisant. Notre impl\'ementation utilise l'API
\texttt{GraphMatcher.isomorphisms\_iter()} de NetworkX.

\subsubsection{G\'en\'erateurs du groupe}

Plut\^ot que d'\'enum\'erer tout $\Aut(\GD)$ (potentiellement
exponentiel en $h$), il suffit de calculer un \emph{ensemble
g\'en\'erateur} $G \subseteq \Aut(\GD)$ : un sous-ensemble dont les
compositions engendrent tout le groupe. Le th\'eor\`eme de Cayley
garantit qu'un tel ensemble de taille $\bigO(\log |\Aut(\GD)|)$
existe ; nous l'extrayons par une proc\'edure gloutonne classique.

\subsection{Th\'eor\`eme de caract\'erisation}

\begin{theoreme}[\'Equivalence redondance--automorphisme]
\label{thm:sym}
Soient $\alpha, \beta : \VD \to \{5, 6, 7\}$ deux solutions du CSP. Les
mol\'ecules associ\'ees \`a $\alpha$ et $\beta$ sont \emph{isomorphes en
tant que graphes de carbone} si et seulement s'il existe $\pi \in
\Aut(\GD)$ tel que $\beta = \pi \cdot \alpha$.
\end{theoreme}

\begin{proof}
\textbf{(}$\Leftarrow$\textbf{)} Soit $\pi \in \Aut(\GD)$ avec
$\beta = \pi \cdot \alpha$. L'automorphisme $\pi$ pr\'eserve l'adjacence
dans $\GD$, donc envoie chaque face $v$ d'$\GB$ sur la face $\pi(v)$ ;
de plus, $\beta(\pi(v)) = \alpha(v)$, donc les tailles sont pr\'eserv\'ees.
Les ar\^etes partag\'ees entre faces adjacentes le sont \'egalement.
On en d\'eduit que $\pi$ induit un isomorphisme entre $G_\alpha$ et
$G_\beta$ (les graphes de carbone respectifs).

\textbf{(}$\Rightarrow$\textbf{)} Supposons $G_\alpha \cong G_\beta$
via un isomorphisme $\varphi : V(G_\alpha) \to V(G_\beta)$.

\begin{enumerate}[label=(\arabic*),leftmargin=*]
  \item \textbf{2-connexit\'e.}
        Dans un benz\'eno\"ide \emph{et toute structure substitu\'ee},
        deux cycles adjacents partagent exactement une ar\^ete
        (= 2~atomes). Aucun sommet n'est point d'articulation : le
        graphe de carbone est 2-connexe.
  \item \textbf{Unicit\'e du plongement planaire (Whitney).}
        Le th\'eor\`eme de Whitney (1933) affirme que tout graphe
        planaire 2-connexe admet un \emph{unique plongement planaire
        combinatoire} (\`a r\'eflexion pr\`es). Cela signifie que
        $\varphi$ pr\'eserve (ou inverse) l'ordre cyclique des ar\^etes
        autour de chaque sommet.
  \item \textbf{Pr\'eservation des faces.}
        L'unicit\'e du plongement implique que $\varphi$ envoie les
        faces de $G_\alpha$ sur les faces de $G_\beta$. Les faces
        internes \'etant exactement les cycles de la mol\'ecule, $\varphi$
        induit une bijection $\pi : \VD \to \VD$.
  \item \textbf{$\pi \in \Aut(\GD)$.}
        Deux faces sont adjacentes dans $\GD$ ssi elles partagent une
        ar\^ete dans le graphe de carbone. Comme $\varphi$ pr\'eserve
        l'adjacence des ar\^etes, $\pi$ pr\'eserve l'adjacence du dual.
  \item \textbf{$\beta = \pi \cdot \alpha$.}
        La taille d'une face est un invariant d'isomorphisme : la face
        $v$ de $G_\alpha$ (de taille $\alpha(v)$) est envoy\'ee sur la face
        $\pi(v)$ de $G_\beta$ (de m\^eme taille). Donc
        $\beta(\pi(v)) = \alpha(v)$, c'est-\`a-dire $\beta = \pi \cdot
        \alpha$.
        \qedhere
\end{enumerate}
\end{proof}

\begin{corollaire}
Les contraintes lex-leader (C4) sur $\Aut(\GD)$ sont \emph{exactes} :
elles \'eliminent toutes les solutions redondantes et seulement elles.
Aucune structure mol\'eculairement distincte n'est perdue.
\end{corollaire}

\subsection{Contraintes lex-leader}

\begin{definition}[Lex-leader]
Pour un ordre fix\'e $(v_1, v_2, \ldots, v_h)$ des sommets et un
automorphisme $\pi \in \Aut(\GD)$, la contrainte lex-leader associ\'ee
\`a $\pi$ est :
\[
  (x_{v_1}, x_{v_2}, \ldots, x_{v_h})
  \;\leq_{\mathrm{lex}}\;
  (x_{\pi(v_1)}, x_{\pi(v_2)}, \ldots, x_{\pi(v_h)}).
\]
\end{definition}

Une solution est dite \emph{lex-leader de son orbite} si elle satisfait
cette in\'egalit\'e pour \emph{tous} les $\pi \in \Aut(\GD)$. Il suffit
en pratique d'imposer la contrainte pour un \emph{ensemble g\'en\'erateur}
de $\Aut(\GD)$, ce qui r\'eduit consid\'erablement le nombre de
contraintes ajout\'ees au mod\`ele.

\begin{exemple}
Pour \texttt{first.graph} ($h = 3$, dual = chemin $v_0 - v_1 - v_2$),
l'unique automorphisme non identit\'e est la r\'eflexion
$\pi : v_0 \leftrightarrow v_2,\; v_1 \mapsto v_1$. La contrainte
lex-leader est :
\[
  (x_{v_0}, x_{v_1}, x_{v_2}) \;\leq_{\mathrm{lex}}\;
  (x_{v_2}, x_{v_1}, x_{v_0}).
\]
La solution $(5, 6, 7)$ est accept\'ee ; son sym\'etrique $(7, 6, 5)$ est
rejet\'e. Tableau des paires d'orbites :
\begin{center}
\renewcommand{\arraystretch}{1.2}
\begin{tabular}{ccc}
\toprule
Orbite & Lex-leader & Rejet\'ee \\
\midrule
$\{(5,6,7), (7,6,5)\}$ & $(5,6,7)$ & $(7,6,5)$ \\
$\{(6,6,6)\}$ & $(6,6,6)$ & --- (orbite singleton) \\
\bottomrule
\end{tabular}
\end{center}
\end{exemple}

\subsection{Justification du traitement en amont}

Le solveur ACE op\`ere en \og boite noire \fg{} : nous lui transmettons
un mod\`ele complet (variables, domaines, contraintes) au format XCSP3
et il retourne les solutions. Il n'y a pas d'acc\`es \`a l'\'etat interne
pendant la r\'esolution (pas de callback, pas d'injection de
\emph{nogoods} \`a la vol\'ee).

C'est pourquoi la rupture de sym\'etrie doit \^etre enti\`erement encod\'ee
\emph{dans le mod\`ele}, avant la r\'esolution :
\begin{enumerate}[leftmargin=*]
  \item Calculer $\Aut(\GD)$ et ses g\'en\'erateurs en Python (NetworkX).
  \item Pour chaque g\'en\'erateur $\pi$, ajouter une contrainte
        \texttt{LexIncreasing} au mod\`ele PyCSP3.
  \item Transmettre le mod\`ele complet au solveur.
\end{enumerate}

Par le Th\'eor\`eme~\ref{thm:sym}, cette approche garantit l'obtention
de toutes les structures distinctes, sans doublon.

% ════════════════════════════════════════════════════════════
\section{Algorithmes de graphe pour la reconstruction 3D}
\label{sec:reconstruction}
% ════════════════════════════════════════════════════════════

Le solveur CSP retourne des assignations $\alpha : \VD \to \{5, 6, 7\}$.
Pour les confronter \`a la validation chimique (xTB) et au test de
planarit\'e, nous devons \textbf{mat\'erialiser} chaque solution en une
g\'eom\'etrie 3D explicite : c'est la phase de \emph{reconstruction}.

Cette phase repose sur trois algorithmes de graphe coupl\'es :
\begin{enumerate}[leftmargin=*]
  \item modification topologique des cycles (contraction d'un sommet
        pour pentagone, insertion d'un sommet pour heptagone) ;
  \item parcours en largeur (BFS) sur le graphe dual pour placer les
        cycles \og un par un \fg{} dans le plan ;
  \item couplage maximum (Blossom) pour l'attribution des liaisons
        doubles ($\pi$-syst\`eme).
\end{enumerate}

Nous d\'ecrivons ici les deux premiers, qui rel\`event de la
combinatoire des graphes ; le troisi\`eme est une application classique
de la th\'eorie du matching et est mentionn\'e \`a la
section~\ref{sec:pipeline}.

\subsection{Phase A : modifications topologiques}

\begin{definition}[Op\'eration de contraction]
\label{def:contraction}
Soit $C$ un cycle (= une face) du graphe de carbone, et $w \in V(C)$
un sommet du cycle. La \emph{contraction} de $w$ dans $C$ identifie
$w$ avec l'un de ses voisins sur $C$, fusionnant les deux ar\^etes
incidentes en une seule. La longueur du cycle d\'ecroit de~1.
\end{definition}

\begin{definition}[Op\'eration d'insertion]
\label{def:insertion}
Soit $C$ un cycle et $e \in E(C)$ une ar\^ete du cycle. L'\emph{insertion}
sur $e$ subdivise $e$ en deux ar\^etes via un nouveau sommet. La longueur
du cycle augmente de~1.
\end{definition}

\subsubsection{Conditions de l\'egalit\'e d'une op\'eration}

L'invariant clef de la phase topologique est : \emph{les ar\^etes
partag\'ees entre deux cycles sont inviolables.}

\begin{lemme}[L\'egalit\'e d'une contraction]
\label{lem:contraction}
Une contraction du sommet $w$ dans le cycle $v$ est \emph{l\'egale}
ssi les deux ar\^etes incidentes \`a $w$ dans $v$ sont libres
(non partag\'ees). En notations : $w \in I(v)$.
\end{lemme}

\begin{lemme}[L\'egalit\'e d'une insertion]
\label{lem:insertion}
Une insertion sur l'ar\^ete $e$ du cycle $v$ est l\'egale ssi $e$ est
libre (non partag\'ee) : $e$ correspond \`a une position $i$ avec
$\sigma_i = 0$.
\end{lemme}

Ces lemmes formalisent les conditions C2bis (Propri\'et\'e~\ref{prop:c2bis})
au niveau de la reconstruction : ce sont les m\^emes pr\'edicats
$I(v) \neq \emptyset$ et $F_v \neq \emptyset$.

\subsubsection{Algorithme de la phase A}

\begin{algorithm}[!ht]
\caption{Phase A : modifications topologiques par cycle.}
\label{alg:topo}
\begin{algorithmic}[1]
\Require Graphe de carbone $\GB$, assignation $\alpha : \VD \to \{5,6,7\}$,
patterns $\sigma(\cdot)$.
\Ensure Graphe modifi\'e $\GB'$ avec cycles de tailles $\alpha(v)$.
\State $\GB' \gets \GB$ \Comment{copie modifiable}
\For{$v \in \VD$}
  \If{$\alpha(v) = 6$}
    \State \textbf{continuer} \Comment{cycle inchang\'e}
  \ElsIf{$\alpha(v) = 5$}
    \State Choisir $w \in I(v)$
    \State Contracter $w$ dans le cycle $v$ de $\GB'$
  \ElsIf{$\alpha(v) = 7$}
    \State Choisir $i$ avec $\sigma_i = 0$
    \State Ins\'erer un nouveau sommet sur le c\^ot\'e $i$ de $v$
  \EndIf
\EndFor
\State \Return $\GB'$
\end{algorithmic}
\end{algorithm}

\begin{propriete}[Terminaison et invariant]
L'algorithme~\ref{alg:topo} termine en $\bigO(h)$ op\'erations
\'el\'ementaires, et maintient l'invariant : toute ar\^ete partag\'ee
$e = (u, v) \in \ED$ reste pr\'esente apr\`es modification, avec ses
deux cycles incidents toujours adjacents par $e$.
\end{propriete}

\subsection{Phase B : placement g\'eom\'etrique par parcours BFS}

Une fois la topologie modifi\'ee, il faut placer chaque cycle dans le
plan euclidien comme un \textbf{polygone r\'egulier}, en respectant les
incidences via les ar\^etes partag\'ees. Notre approche est un
\textbf{parcours en largeur} (BFS) sur $\GD$.

\subsubsection{Polygones r\'eguliers : g\'eom\'etrie de base}

\begin{definition}[Polygone r\'egulier]
Le polygone r\'egulier \`a $n$ c\^ot\'es de longueur $a$ a pour
\emph{rayon} (distance centre--sommet) et \emph{apoth\`eme} (distance
centre--milieu de c\^ot\'e) :
\[
  R(n) \;=\; \frac{a}{2 \sin(\pi/n)},
  \qquad
  r(n) \;=\; \frac{a}{2 \tan(\pi/n)}.
\]
\end{definition}

Avec $a = 1{,}40$~\AA{} (longueur C--C aromatique typique), nous avons :
\begin{center}
\begin{tabular}{cccc}
\toprule
$n$ & 5 & 6 & 7 \\
\midrule
$R(n)$ (\AA) & 1{,}19 & 1{,}40 & 1{,}61 \\
$r(n)$ (\AA) & 0{,}96 & 1{,}21 & 1{,}45 \\
\bottomrule
\end{tabular}
\end{center}

\subsubsection{Algorithme BFS de placement}

L'id\'ee est de :
\begin{enumerate}[leftmargin=*]
  \item choisir une \emph{racine} $r \in \VD$ (par exemple le sommet
        de plus haut degr\'e, pour ancrer la mol\'ecule sur la zone la
        plus contrainte) ;
  \item placer le cycle $r$ comme polygone r\'egulier centr\'e \`a
        l'origine ;
  \item parcourir $\GD$ en BFS \`a partir de $r$ : pour chaque sommet
        enfant $c$ d\'ecouvert via une ar\^ete partag\'ee $e = (p, c)$
        avec son parent $p$, placer le cycle $c$ en \og collant \fg{}
        son c\^ot\'e $\posIn{c}{e}$ sur le c\^ot\'e $\posIn{p}{e}$ de $p$.
\end{enumerate}

\begin{algorithm}[!ht]
\caption{Phase B : placement g\'eom\'etrique par BFS sur $\GD$.}
\label{alg:bfs-placement}
\begin{algorithmic}[1]
\Require Graphe dual $\GD$, assignation $\alpha$, positions $\posIn{\cdot}{\cdot}$.
\Ensure Coordonn\'ees 3D pour chaque atome (z = 0).
\State $r \gets \arg\max_{v \in \VD} \degv$ \Comment{racine = sommet de degr\'e max}
\State Placer le cycle $r$ comme polygone r\'egulier centr\'e \`a l'origine, orient\'e $\theta_r = 0$
\State $Q \gets$ file FIFO contenant $r$
\State $\text{visit\'e} \gets \{r\}$
\While{$Q \neq \emptyset$}
  \State $p \gets Q.\text{d\'efiler}()$
  \For{$c \in \NN(p)$ avec $c \notin \text{visit\'e}$}
    \State $e \gets (p, c)$
    \State Soit $M$ le milieu du c\^ot\'e $\posIn{p}{e}$ de $p$ (d\'ej\`a plac\'e)
    \State $\theta_c \gets$ orientation telle que le c\^ot\'e $\posIn{c}{e}$ co\"incide avec $M$
    \State Placer le cycle $c$ comme polygone r\'egulier de taille $\alpha(c)$, centre \`a apoth\`eme $r(\alpha(c))$ de $M$, orient\'e $\theta_c$
    \State $\text{visit\'e} \gets \text{visit\'e} \cup \{c\}$
    \State $Q.\text{enfiler}(c)$
  \EndFor
\EndWhile
\State \Return Coordonn\'ees de tous les atomes
\end{algorithmic}
\end{algorithm}

\subsubsection{Coh\'erence du parcours BFS}

\begin{lemme}[Existence d'un placement]
\label{lem:bfs-existence}
Soit $\alpha$ une solution du CSP. L'algorithme~\ref{alg:bfs-placement}
termine et produit un plongement g\'eom\'etrique \emph{valide} (sans
recouvrement des cycles) ssi la condition C2bis
(Propri\'et\'e~\ref{prop:c2bis}) est satisfaite \emph{et} la configuration
locale autour de chaque sommet appartient \`a la table $\Tn$
(contrainte~C3).
\end{lemme}

\begin{proof}[Esquisse]
Le BFS visite chaque sommet exactement une fois ($\bigO(h)$). Pour chaque
sommet plac\'e, la coh\'erence g\'eom\'etrique repose sur :
\begin{itemize}
  \item la longueur des ar\^etes : $a = 1{,}40$~\AA{} pour tous les
        polygones (h\'eritage du r\'eseau aromatique) ;
  \item l'orientation : impos\'ee par le placement du parent et la
        position $\posIn{c}{e}$.
\end{itemize}
Si tous les voisinages appartiennent \`a $\Tn$, alors localement la
configuration est connue plane (par construction de la table). La somme
des d\'eviations angulaires autour de chaque sommet est nulle (formule
d'Euler discr\'etis\'ee). En l'absence d'incompatibilit\'es globales, le
plongement BFS produit un plongement valide.
\end{proof}

\begin{remarque}[Cas d'\'echec g\'eom\'etrique]
M\^eme avec C2bis et C3 satisfaites, il peut subsister des solutions
\og localement valides mais globalement infaisables \fg{} sur les
benz\'eno\"ides \`a fortes contraintes (h9 en mode
\texttt{--no-freeze --no-table}). Ces \'echecs proviennent de
discr\'epances angulaires cumul\'ees sur des \emph{cycles longs} de
$\GD$, et constituent un probl\`eme ouvert (cf. courbure de Gauss
discr\`ete sur surfaces polygonales).
\end{remarque}

\subsubsection{Illustration : un cas \`a 4 hexagones substitu\'e}

\begin{figure}[!ht]
\centering
\begin{tikzpicture}[scale=0.7, every node/.style={font=\small}]
  % Pentagone (rouge clair)
  \begin{scope}[shift={(0,0)}]
  \foreach \i in {0,...,4} {
    \coordinate (P\i) at ({cos(72*\i+90)*1.19}, {sin(72*\i+90)*1.19});
  }
  \draw[thick,fill=jaunePent] (P0) -- (P1) -- (P2) -- (P3) -- (P4) -- cycle;
  \node at (0,0) {\small 5};
  \foreach \i in {0,...,4} {\fill (P\i) circle (1.5pt);}
  \end{scope}
  % Hexagone (gris)
  \begin{scope}[shift={(2.42,0)}]
  \foreach \i in {0,...,5} {
    \coordinate (H\i) at ({cos(60*\i+30)*1.40}, {sin(60*\i+30)*1.40});
  }
  \draw[thick,fill=grisHex] (H0) -- (H1) -- (H2) -- (H3) -- (H4) -- (H5) -- cycle;
  \node at (0,0) {\small 6};
  \foreach \i in {0,...,5} {\fill (H\i) circle (1.5pt);}
  \end{scope}
  % Heptagone (indigo)
  \begin{scope}[shift={(5.30,0)}]
  \foreach \i in {0,...,6} {
    \coordinate (G\i) at ({cos(51.43*\i+90)*1.61}, {sin(51.43*\i+90)*1.61});
  }
  \draw[thick,fill=indigoHept] (G0) -- (G1) -- (G2) -- (G3) -- (G4) -- (G5) -- (G6) -- cycle;
  \node at (0,0) {\small 7};
  \foreach \i in {0,...,6} {\fill (G\i) circle (1.5pt);}
  \end{scope}
\end{tikzpicture}
\caption{Trois polygones r\'eguliers (pentagone, hexagone, heptagone) de
m\^eme longueur de c\^ot\'e $a = 1{,}40$~\AA{}, fusionn\'es par leurs
c\^ot\'es partag\'es. Les rayons $R(n)$ diff\`erent (1{,}19 / 1{,}40 /
1{,}61~\AA{}), induisant une distorsion locale absorb\'ee par le
voisinage.}
\label{fig:trois-polygones}
\end{figure}

\subsection{Le graphe dual r\'esolu : graphe de placement}

Une fois le BFS termin\'e, chaque sommet $v$ de $\GD$ est annot\'e par :
\begin{itemize}
  \item ses coordonn\'ees cart\'esiennes 2D (centre $c_v$ et orientation
        $\theta_v$) ;
  \item la liste de ses sommets de polygone (= positions des atomes de
        carbone).
\end{itemize}

L'ensemble de ces annotations forme le \emph{graphe de placement} : une
structure interm\'ediaire entre le mod\`ele topologique (le dual) et la
g\'eom\'etrie atomique (le graphe de carbone enrichi par les
coordonn\'ees).

\subsection{Synth\`ese : pipeline reconstruction}

\begin{figure}[!ht]
\centering
\begin{tikzpicture}[
  node distance=12mm,
  every node/.style={font=\small},
  box/.style={draw,thick,rounded corners=2pt,fill=bleuSub!10,minimum width=27mm,minimum height=10mm,align=center},
  arr/.style={->,thick,>=Stealth}
]
  \node[box] (csp) {Solution CSP\\$\alpha : \VD \to \{5,6,7\}$};
  \node[box, right=of csp] (topo) {Phase A\\Modifications\\topologiques};
  \node[box, right=of topo] (bfs) {Phase B\\BFS placement\\($\GD$)};
  \node[box, below=of topo] (assemble) {Phase C\\Assemblage\\Blossom matching};
  \node[box, right=of assemble] (xyz) {Fichier XYZ\\(3D + H)};

  \draw[arr] (csp) -- (topo);
  \draw[arr] (topo) -- (bfs);
  \draw[arr] (bfs) -- ($(assemble.north)+(0,0)$);
  \draw[arr] (assemble) -- (xyz);
\end{tikzpicture}
\caption{Pipeline de reconstruction 3D depuis une solution CSP. Les
phases A et B (topologie + BFS) sont d\'ecrites en
section~\ref{sec:reconstruction} ; la phase C (couplage maximum pour
liaisons doubles, ajout des hydrog\`enes) est trait\'ee dans le rapport
principal.}
\label{fig:reconstruction-pipeline}
\end{figure}

% ════════════════════════════════════════════════════════════
\section{Pipeline exp\'erimental et complexit\'e globale}
\label{sec:pipeline}
% ════════════════════════════════════════════════════════════

\subsection{Vue d'ensemble du pipeline}

Le pipeline complet, depuis un fichier \texttt{.graph} Benzai jusqu'aux
g\'eom\'etries 3D valid\'ees, comprend les \'etapes suivantes :

\begin{enumerate}[leftmargin=*]
  \item \textbf{Parsing} (\texttt{parser.py}) : lecture du fichier
        Benzai, construction de $\GB$ et $\GD$, calcul de
        $\sigma(\cdot)$, $b(\cdot)$, $I(\cdot)$, $F_\cdot$.
        Complexit\'e : $\bigO(|V_\B| + |E_\B| + h^2)$.
  \item \textbf{Pr\'e-traitement CSP} (\texttt{preprocessing.py}) :
        application de C2 et C2bis, filtrage des tables, calcul des
        automorphismes. Complexit\'e : $\bigO(h^2 + h \cdot |\Tn|)$ pour
        le filtrage ; $\bigO(\text{poly}(h))$ en pratique pour
        $\Aut(\GD)$ via VF2.
  \item \textbf{Compilation XCSP3} (\texttt{model.py}) : g\'en\'eration
        du fichier \texttt{model.xml} via PyCSP3. Lin\'eaire en taille
        du mod\`ele.
  \item \textbf{R\'esolution} (ACE) : appel\'e en sous-processus, ACE
        \'enum\`ere toutes les solutions satisfaisant C1+C3+C4. Les
        contraintes C2/C2bis sont d\'ej\`a int\'egr\'ees comme
        r\'eductions de domaine.
  \item \textbf{Reconstruction 3D} (\texttt{reconstruction/}) : phases
        A, B, C de la section~\ref{sec:reconstruction}.
  \item \textbf{Validation xTB + planarit\'e} : optimisation
        g\'eom\'etrique semi-empirique et test ACP des plans de cycles.
\end{enumerate}

\subsection{Complexit\'es \'etape par \'etape}

\begin{table}[!ht]
\centering
\renewcommand{\arraystretch}{1.3}
\begin{tabular}{@{}llp{6.5cm}@{}}
\toprule
\textbf{\'Etape} & \textbf{Complexit\'e} & \textbf{Comm.} \\
\midrule
Parsing Benzai & $\bigO(|V_\B| + |E_\B|)$ & lecture lin\'eaire \\
Construction $\GD$ & $\bigO(h^2)$ & double boucle sur hexagones \\
Calcul $\sigma, b, I, F$ & $\bigO(h)$ & temps constant par sommet \\
Pr\'e-traitement C2/C2bis & $\bigO(h)$ & r\'eduction de domaines \\
Filtrage tables & $\bigO(h \cdot |\Tn|_{\max})$ & filtrage par pattern \\
Arc-consistance & $\bigO(|\ED| \cdot |\Tn|_{\max})$ & propagation \\
Calcul $\Aut(\GD)$ & $\bigO(\text{poly}(h))$ & VF2 (quasi-poly th\'eorique) \\
R\'esolution ACE & $\bigO(3^h)$ pire cas & exponentielle dans $h$, ex\'ecution sous-seconde en pratique pour $h \leq 6$ \\
Reconstruction (par sol) & $\bigO(h)$ & topologie + BFS \\
Validation xTB & $\bigO(N^3)$--$\bigO(N^4)$ & $N$ = nombre d'atomes ; dominant en temps r\'eel \\
\bottomrule
\end{tabular}
\caption{Complexit\'es des \'etapes du pipeline. Toutes les \'etapes
hors r\'esolution CSP et validation xTB sont polynomiales en $h$.}
\label{tab:complexites}
\end{table}

\subsection{R\'esultats exp\'erimentaux}

Les benz\'eno\"ides de notre jeu de tests vont de $h = 2$ \`a $h = 20$
hexagones. Les performances mesur\'ees :
\begin{itemize}
  \item \textbf{Pr\'e-traitement complet} : $< 100$~ms pour $h \leq 20$ ;
  \item \textbf{R\'esolution ACE} : $< 1$~seconde pour $h \leq 16$
        (\texttt{second.graph}, 41 solutions \'enum\'er\'ees) ;
  \item \textbf{Reconstruction + validation xTB} : 10--60~secondes
        par solution selon la taille (\'etape dominante).
\end{itemize}

Pour des benz\'eno\"ides plus grands ($h_8$, $h_9$), la r\'esolution
reste sous-seconde gr\^ace au pr\'e-traitement agressif, mais le
volume de solutions \`a valider explose (jusqu'\`a $\sim 1{,}5 \cdot 10^6$
sur $h_9$), n\'ecessitant un d\'eploiement sur cluster.

\subsection{Comparaison aux approches alternatives}

\subsubsection{\'Enum\'eration na\"ive}

Sur l'espace brut $\{5,6,7\}^h$, l'\'enum\'eration na\"ive teste $3^h$
configurations. M\^eme avec un test rapide (C1 + C2 en $\bigO(h)$), pour
$h = 20$ cela repr\'esente $\sim 3{,}5 \cdot 10^9$ tests, soit
$\sim 1$ heure \`a $10^6$ tests/seconde. Notre approche CSP r\'esout en
$< 1$ seconde gr\^ace au pr\'e-traitement et \`a la propagation.

\subsubsection{Approche de Varet (2022)}

La th\`ese de Varet utilise Choco-Solver (Java) avec injection
dynamique de \emph{nogoods} pour la rupture de sym\'etrie. Notre
approche par lex-leader est plus simple \`a impl\'ementer (compatible
\og boite noire \fg{}) et offre les m\^emes garanties th\'eoriques
(Th\'eor\`eme~\ref{thm:sym}). En contrepartie, le solveur ne peut pas
exploiter les nogoods dynamiques pour acc\'el\'erer le retour arri\`ere ;
mais la diff\'erence est marginale sur nos tailles.

% ════════════════════════════════════════════════════════════
\section{Conclusion du chapitre}
\label{sec:conclusion}
% ════════════════════════════════════════════════════════════

Nous avons pr\'esent\'e dans ce chapitre une mod\'elisation compl\`ete
du probl\`eme de g\'en\'eration de mol\'ecules non-benz\'eno\"ides comme
un probl\`eme de satisfaction de contraintes sur le graphe dual du
benz\'eno\"ide d'entr\'ee. Les contributions th\'eoriques principales
sont :
\begin{enumerate}[leftmargin=*]
  \item Une caract\'erisation rigoureuse des \emph{hexagones immobiles}
        en termes du seul pattern $\sigma(v)$
        (Propri\'et\'es~\ref{prop:degre6} et~\ref{prop:b2}), permettant
        une r\'eduction de domaine en $\bigO(h)$.
  \item L'introduction d'une condition n\'ecessaire suppl\'ementaire
        de r\'ealisabilit\'e topologique
        (Propri\'et\'e~\ref{prop:c2bis} : sommets int\'erieurs libres
        pour pentagones, c\^ot\'es libres pour heptagones), qui filtre
        jusqu'\`a 75~\% des solutions infaisables sur les benz\'eno\"ides
        d'ordre $\geq 8$.
  \item La preuve compl\`ete de la conservation du carbone via la
        formule d'Euler appliqu\'ee au graphe de carbone
        (Propri\'et\'e~\ref{prop:euler}).
  \item La preuve que les automorphismes du graphe dual \emph{captent
        exactement} les redondances entre solutions, en s'appuyant sur
        la 2-connexit\'e des benz\'eno\"ides et le th\'eor\`eme
        d'unicit\'e du plongement planaire de Whitney
        (Th\'eor\`eme~\ref{thm:sym}).
  \item Un parcours BFS sur le graphe dual pour la reconstruction
        g\'eom\'etrique des solutions, avec preuve d'existence d'un
        placement valide (Lemme~\ref{lem:bfs-existence}).
\end{enumerate}

Du point de vue algorithmique, le pipeline construit allie th\'eorie
des graphes (automorphismes, parcours BFS, propri\'et\'es structurelles
du dual), th\'eorie des CSP (contraintes extensionnelles,
arc-consistance, lex-leader) et th\'eorie du matching (couplage maximum
pour les liaisons doubles, mentionn\'e en \ref{sec:pipeline}). Le tout
est polynomial en $h$ \`a l'exception du c\oe ur de la r\'esolution CSP
(exponentiel dans le pire cas, mais sub-seconde en pratique gr\^ace au
pr\'e-traitement).

\paragraph{Perspectives th\'eoriques.}
Plusieurs questions restent ouvertes :
\begin{itemize}
  \item \textbf{Contraintes globales de planarit\'e} : nos contraintes
        actuelles sont \emph{locales} (par sommet + voisinage imm\'ediat).
        Pour les grands benz\'eno\"ides, un \og d\'eficit angulaire
        cumul\'e \fg{} peut emp\^echer la planarit\'e globale, sans
        \^etre d\'etect\'e par C2/C2bis/C3. Une formulation en termes
        de \emph{courbure de Gauss discr\`ete} (somme des d\'eficits
        angulaires) reste \`a explorer.
  \item \textbf{Sym\'etries dynamiques} : la rupture lex-leader est
        statique. Des sym\'etries qui n'apparaissent qu'apr\`es certaines
        affectations (\emph{sym\'etries conditionnelles}) ne sont pas
        captur\'ees ; cela pourrait gagner un facteur sur les tr\`es
        grands benz\'eno\"ides.
  \item \textbf{D\'enombrement direct} : peut-on calculer le nombre
        de solutions \emph{sans} les \'enum\'erer (fonction
        g\'en\'eratrice, comptage par orbites \`a la Burnside) ? Cela
        permettrait des \'etudes statistiques sur les classes de
        benz\'eno\"ides.
\end{itemize}

Le chapitre suivant aborde la phase \emph{aval} du pipeline : validation
physico-chimique par xTB, test de planarit\'e par analyse en composantes
principales, et statistiques sur les structures \'enum\'er\'ees.

\end{document}
